\chapter{Introduction}
\ChapterLabel{Introduction}
\index{introduction}
\index{Chapter!Introduction}
\index{normative language}
\index{priority}
\index{source of truth}
\index{V-model}

\section{Chapter Responsibility}
\label{Section:IntroductionResponsibility}
Establish the purpose, scope, audience, and conventions of the engineering record.

\noindent\textbf{Expected result:} A shared understanding of how to read and use the architecture.

\section{How to Use This Chapter}
State the document purpose, scope, audience, conventions, and relationship to
other project artifacts. Replace bracketed text with project-specific content.
Keep binding obligations in the Requirements chapter.

\section{Purpose}
This document is the authoritative engineering specification for
\ProjectName. It records why the system is needed, what it shall do, how it is
structured, how it will be tested, and what evidence is required before
release. It is readable by people and structured for use by engineering and
\gls{ai}-assisted development tools.

\begin{dsnNote}[Central lesson]
A vague request produces unpredictable software. A precise, traceable
specification allows engineers to generate, review, test, and, when necessary,
regenerate a conforming implementation.
\end{dsnNote}

\section{Scope}
Define the system boundary, included capabilities, exclusions, operating
environment, major constraints, and the release covered by this document.

\textit{Generic example: \ProjectName\ enables [primary user] to [achieve
outcome]. This release includes [included capabilities] and excludes
[out-of-scope capabilities]. It operates in [target environment] and depends
on [external systems or services].}

\section{Audience}
Identify the people who create, approve, implement, test, operate, maintain,
audit, or use the system. State which sections are most relevant to each role.

\textit{Generic example: The primary audiences are customer representatives,
systems engineers, architects, developers, testers, operators, maintainers,
security reviewers, and project approvers. One person may perform several
roles on a small project.}

\section{Normative Language and Priority}
Use \textbf{shall} for every binding requirement. Do not use \textbf{should} as
guidance in normative statements. Assign each requirement one priority:
\textbf{\gls{must}}, \textbf{\gls{should}}, \textbf{\gls{could}}, or
\textbf{\gls{would}}. Priority expresses delivery importance; it does not change
whether a statement containing \textbf{shall} is binding for its assigned
implementation version.

\section{Document as Source of Truth}
Requirements alone do not define a reproducible implementation. The complete
source of truth consists of stakeholder needs, user stories, requirements, use
cases, architecture and interface decisions, implementation constraints,
tests, verification evidence, validation evidence, and release criteria.
Generated code is an implementation artifact governed by this specification.

\section{Architecture as Maintained Input for AI Design}
\label{Section:ArchitectureAIDesign}
\index{AI!architecture as maintained input}
\index{architecture!AI-assisted design}
Figure~\ref{Figure:MaintainedSourceComparison} distinguishes two workflows.
A compiler regenerates lower-level output from reviewed C/C++ source and its
build configuration. An AI request is often incomplete: it may omit interface
assumptions, error behavior, or earlier decisions. Generated code therefore
requires review, tests, and version control. Repeating a prompt does not
reproduce an approved implementation by itself.

\begin{figure}[htbp]
\centering
\resizebox{\textwidth}{!}{% Adapted from MSE Figure 2.1.
% Adapted from Week-02/TikZ/sldSourceOfTruthAnalogy/sldSourceOfTruthAnalogyDrawing.tex.
\begin{tikzpicture}[font=\small,
 source/.style={draw=dsnBlue,fill=dsnLightBlue,rounded corners=2pt,
   text width=3.0cm,minimum height=1.25cm,align=center},
 process/.style={draw=black!70,fill=black!6,rounded corners=2pt,
   text width=2.3cm,minimum height=1.25cm,align=center},
 output/.style={draw=black!55,rounded corners=2pt,
   text width=3.7cm,minimum height=1.25cm,align=center},
 arrow/.style={-{Latex[length=2.2mm]},line width=.8pt}]
\node[font=\small\bfseries,anchor=east] at (-1.1,0) {Compiler};
\node[source] (source) at (0.7,0) {C/C++ source\\stored and reviewed};
\node[process] (compiler) at (4.5,0) {Compiler};
\node[output] (binary) at (8.6,0) {Assembly /\\machine code\\regenerated};
\draw[arrow] (source) -- (compiler);
\draw[arrow] (compiler) -- (binary);
\node[font=\small\bfseries,anchor=east] at (-1.1,-2) {AI today};
\node[output,text width=3cm] (prompt) at (.7,-2) {Prompt\\usually incomplete};
\node[process] (ai) at (4.5,-2) {AI};
\node[source,text width=3.7cm] (code) at (8.6,-2) {Generated code\\stored and reviewed};
\draw[arrow] (prompt) -- (ai);
\draw[arrow] (ai) -- (code);
\end{tikzpicture}
}
\caption{Maintained source in compiler and AI workflows, adapted from MSE Figure 2.1.}
\label{Figure:MaintainedSourceComparison}
\end{figure}

The goal of this template is to make the approved architecture the maintained
engineering input to AI-assisted design and implementation. Its linked needs,
requirements, interfaces, responsibilities, constraints, and rationale give a
bounded task its context. The task request identifies the affected baseline and
change; the architecture remains the authoritative record of design intent.
This analogy does not make an AI agent a deterministic compiler. Engineers must
review, verify, and version its code and tests, and resolve missing acceptance
criteria before claiming conformance.

Figure~\ref{Figure:ArchitectureUpdateFeedback} shows how implementation and
test findings can improve the architecture. An engineer may edit the affected
views directly, or an AI agent may prepare a proposed update. In either route,
authorized reviewers assess the evidence and approve changes before they enter
the baseline. A finding may affect requirements, interfaces, allocations, and
planned checks together; update those links before using the revised baseline
for the next implementation task. \ChapterRef{AIAgentGuide} and
\ChapterRef{AIImplementationWorkflow} define the task and execution controls.

\begin{figure}[htbp]
\centering
\resizebox{\textwidth}{!}{% Adapted from MSE Figure 2.2.
% Adapted from Week-02/TikZ/sldArchitectureDirectsAI/sldArchitectureDirectsAIDrawing.tex.
% The two return routes share baseline approval, irrespective of who authors the update.
\begin{tikzpicture}[font=\small,
 source/.style={draw=dsnBlue,fill=dsnLightBlue,rounded corners=2pt,
   text width=3.5cm,minimum height=1.35cm,align=center},
 process/.style={draw=black!70,fill=black!6,rounded corners=2pt,
   text width=3.1cm,minimum height=1.1cm,align=center},
 arrow/.style={-{Latex[length=2.2mm]},line width=.8pt},
 feedback/.style={arrow,draw=dsnBlue}]
\node[source] (architecture) at (0,0)
 {LaTeX architecture\\approved baseline\\in source control};
\node[process] (implementation) at (5,0) {AI-assisted\\implementation};
\node[source] (code) at (10,0) {Code and tests\\reviewed, verified,\\versioned};
\draw[arrow] (architecture) -- (implementation);
\draw[arrow] (implementation) -- (code);
\node[process] (findings) at (10,-2.6) {Design findings\\from code and\\test evidence};
\node[process] (human) at (5,-2.6) {Human\\architecture update};
\node[process] (agent) at (5,-5) {AI-assisted\\architecture update};
\node[source] (review) at (0,-2.6) {Review and approval\\of architecture\\changes};
\draw[feedback] (code) -- (findings);
\draw[feedback] (findings) -- (human);
\draw[feedback] (human) -- (review);
\draw[feedback] (findings.south) |- (agent.east);
\draw[feedback] (agent.west) -| (review.south);
\draw[feedback] (review) -- (architecture);
\end{tikzpicture}
}
\caption{Human and AI-assisted architecture feedback, adapted from MSE Figure 2.2.}
\label{Figure:ArchitectureUpdateFeedback}
\end{figure}

\section{Initializing a New Project from This Template}
\label{Section:InitializeNewProject}
\index{template!project initialization}
\index{AI agent!project initialization}
Create a project-specific architecture book from a copy of this template;
keep the master template generic. Use
\hyperref[Skill:AIS-02]{AIS-02: Initialize a Project} for the authoritative
Context, Steps, and Acceptance, including confirmation of project information,
review of supplied evidence, and replacement of supported placeholders.

To start initialization with an AI agent, use the request in
\hyperref[Section:AIInitializePrompt]{Section~\ref{Section:AIInitializePrompt},
Initialize the Architecture Book}. Execute the task using
\ChapterRef{AIImplementationWorkflow}.

\section{Template Tailoring and Domain Profiles}
\label{Section:DomainProfiles}
\index{domain tailoring profile}
\index{template!tailoring}
The core chapters form a generic software-architecture template. A domain
tailoring profile augments those chapters with concerns, artifacts, diagrams,
and review criteria for a particular class of system. It does not replace a
core chapter or create a second authoritative location for project facts.
Requirements, interfaces, architecture decisions, tests, and evidence remain
in their normal core chapters and link to one another through their stable IDs.

Select applicable profiles during project initialization. More than one
profile may apply. Record each applicability decision and rationale in
\ChapterRef{DevelopmentPlan}, retain applicable profile appendices as
authoritative tailoring guidance, and exclude an inapplicable profile from an
instantiated project only after recording the decision. If profile guidance
conflicts, identify the conflict and obtain an authorized human decision before
editing controlled project content.

\small
\begin{tabularx}{\linewidth}{L{4.0cm}L{3.8cm}Y}
\caption{Domain Profile Applicability\label{Table:DomainProfileApplicability}}\\
\toprule
\textbf{Profile} & \textbf{Applicability} & \textbf{Decision and Rationale}\\
\midrule
\endfirsthead
\multicolumn{3}{l}{\small\itshape Domain Profile Applicability (continued)}\\
\toprule
\textbf{Profile} & \textbf{Applicability} & \textbf{Decision and Rationale}\\
\midrule
\endhead
\bottomrule
\endlastfoot
\hyperref[Chapter:EmbeddedSystemsProfile]{Embedded Systems} &
\textit{[Applicable / Not Applicable / TBD]} &
\textit{[System characteristics, scope, and decision authority]}\\
\textit{[Additional profile]} &
\textit{[Applicable / Not Applicable / TBD]} &
\textit{[Reason and relationship to other profiles]}\\
\end{tabularx}
\normalsize

When the Embedded Systems profile is Applicable or TBD, read
\ChapterRef{EmbeddedSystemsProfile} during initialization and for every task
affected by hardware, firmware, real-time behavior, constrained resources,
physical interfaces, safety, security, manufacturing, or field service.

\section{Using This Template with AI Agents}
\label{Section:UsingTemplateAI}
\index{AI agent!architecture workflow}
Use \ChapterRef{AIAgentGuide} to select AI Skills and read their
linked Context, Steps, and Acceptance. Use
\ChapterRef{AIImplementationWorkflow} for execution with one or multiple
agents. To start a task, use
\hyperref[Section:AIStartSelectedTaskPrompt]{Start Any Selected AIS};
\hyperref[Section:AIExamplePrompts]{the task-specific example requests}
provide additional starting prompts. The
\hyperref[Figure:AIGuideActivity]{guide activity diagram} shows the reading
route from repository bootstrap files to task selection and execution.

\section{V-Model Organization}
Figure~\ref{Figure:VModel} connects development decisions with evidence.
Stakeholder needs motivate requirements, which guide architecture, detailed
design, and implementation. Unit tests address detailed behavior; integration
tests address interactions; system verification assesses specified obligations;
and stakeholder validation assesses intended use. These correspondences help
identify the evidence needed for each design decision.

The model includes iteration. Validation findings can return to stakeholder
needs during development. After deployment, operational findings and change
requests can trigger another cycle. A changed baseline requires review of its
requirements, design links, and affected checks. This organization is
compatible with the system and software life-cycle process frameworks in
\cite{ISO15288,ISO12207}; those standards do not prescribe one life-cycle model.

\begin{figure}[htbp]
\centering
\resizebox{\textwidth}{!}{% Adapted from MSE Figure 1.2; keep lifecycle and feedback routes orthogonal.
\begingroup
\begin{tikzpicture}[yscale=1.12,
  stage/.style={draw=dsnBlue,fill=dsnBlue!7,rounded corners=2pt,text width=3.7cm,inner sep=0.10cm,minimum width=3.9cm,minimum height=1.25cm,align=center,font=\fontsize{16}{18}\selectfont},
  evidence/.style={stage,draw=dsnRed,fill=dsnRed!7},
  implementation/.style={stage,draw=dsnGray,fill=dsnGray!10},
  operation/.style={stage,draw=dsnGray,fill=dsnGray!10},
  lifecycle/.style={-{Latex[length=2.2mm]},line width=0.75pt},
  relation/.style={draw=dsnGray,dashed,line width=0.6pt},
  feedback/.style={-{Latex[length=2.2mm]},draw=dsnGray,dashed,line width=0.7pt},
  label/.style={fill=white,inner sep=1.2pt,font=\fontsize{14}{16}\selectfont,text=dsnGray}
]
  \node[stage] (needs) at (-9.2,8.4) {Stakeholder needs};
  \node[stage] (requirements) at (-6.9,6.3) {System and software\\requirements};
  \node[stage] (architecture) at (-4.6,4.2) {Architecture};
  \node[stage] (design) at (-2.3,2.1) {Detailed design};
  \node[implementation] (implementation) at (0,0) {Implementation};
  \node[evidence] (unit) at (2.3,2.1) {Unit tests};
  \node[evidence] (integration) at (4.6,4.2) {Integration tests};
  \node[evidence] (verification) at (6.9,6.3) {System\\verification};
  \node[evidence] (validation) at (9.2,8.4) {Stakeholder\\validation};
  \node[operation] (deployment) at (11.5,10.5) {Deployment};
  \node[operation] (maintenance) at (13.8,12.6) {Operations \&\\ Maintenance};
  \draw[lifecycle,draw=dsnBlue] (needs.south) -- ++(0,-0.4) -| (requirements.north);
  \draw[lifecycle,draw=dsnBlue] (requirements.south) -- ++(0,-0.4) -| (architecture.north);
  \draw[lifecycle,draw=dsnBlue] (architecture.south) -- ++(0,-0.4) -| (design.north);
  \draw[lifecycle,draw=dsnBlue] (design.south) -- ++(0,-0.4) -| ([xshift=-3mm]implementation.north);
  \draw[lifecycle,draw=dsnRed] ([xshift=3mm]implementation.north) -- ++(0,0.4) -| (unit.south);
  \draw[lifecycle,draw=dsnRed] (unit.north) -- ++(0,0.4) -| (integration.south);
  \draw[lifecycle,draw=dsnRed] (integration.north) -- ++(0,0.4) -| (verification.south);
  \draw[lifecycle,draw=dsnRed] (verification.north) -- ++(0,0.4) -| (validation.south);
  \draw[lifecycle,draw=dsnGray] (validation.north) -- ++(0,0.4) -| (deployment.south);
  \draw[lifecycle,draw=dsnGray] (deployment.north) -- ++(0,0.4) -| (maintenance.south);
  \draw[feedback] (maintenance.north) -- (13.8,14.1)
    -- node[label,above=1pt,pos=0.48]{operational feedback / change request} (-9.2,14.1)
    -- (needs.north);
  \draw[relation] (needs) -- node[label,above=1pt]{acceptance} (validation);
  \draw[relation] (requirements) -- node[label,above=1pt]{compliance} (verification);
  \draw[relation] (architecture) -- node[label,above=1pt]{integration} (integration);
  \draw[relation] (design) -- node[label,above=20pt]{correctness} (unit);
  \draw[feedback,draw=dsnBlue] ([xshift=-4mm]validation.north)
    -- (8.8,11.0)
    -- node[label,above=1pt,text=dsnBlue]{agile iteration} (-8.8,11.0)
    -- ([xshift=4mm]needs.north);
\end{tikzpicture}
\endgroup
}
\caption{V-model relationships, agile iteration, and operational feedback, adapted from MSE Figure 1.2.}
\label{Figure:VModel}
\end{figure}
